\documentclass[11pt]{article}
\usepackage[a4paper,margin=1in]{geometry}
\usepackage{amsmath,amssymb,amsthm}
\usepackage{booktabs}
\usepackage{longtable}
\usepackage{xcolor}
\usepackage{hyperref}
\hypersetup{colorlinks=true,linkcolor=blue,urlcolor=blue,citecolor=blue}
\newcommand{\fn}[1]{\texttt{#1}}
\theoremstyle{plain}
\newtheorem{proposition}{Proposition}
\newtheorem{theorem}{Theorem}
\theoremstyle{remark}
\newtheorem{remark}{Remark}

\title{\textbf{Report R27 --- the evolution operator's Schmidt rank}\\
\large What the literature establishes, and why a linear-dimensional module
forces linear growth\\[2pt]
\normalsize rev.\ 2: the bound is now two-sided, and the law is exact}
\author{Tier 1 analysis}
\date{2026-08-23}

\begin{document}
\maketitle
\sloppy
\emergencystretch=3em

\begin{abstract}
Rule $156$ has an operator Schmidt rank $\chi(t)$ of its \emph{evolution
operator} that grows linearly in $t$, in a clean, disorder-free, translation-%
invariant Floquet circuit. This report establishes what the literature says
about that observable and supplies a mechanism. \S\ref{sec:protocol} is an
audited pass over the twelve references supplied plus eleven found by search,
sorted by relevance, each with the finding that bears on our claim.
\textbf{The baseline is Dubail's}: the OSEE of $U(t)=e^{-iHt}$ grows
\emph{linearly}, hence $\chi$ \emph{exponentially}, ``unless the Hamiltonian is
in a localized phase''. Zhou and Luitz classify the three observed growth laws
for $U(t)$ --- linear, power law, logarithmic --- and both sub-linear classes
occur only with disorder. Every disorder-free sub-linear result we could find,
including Alba--Dubail--Medenjak for Rule~54 and the very recent
Tan--Prosen clean non-integrable circuit, is about \emph{local operators in the
Heisenberg picture}, which Dubail explicitly separates from the propagator. We
found no disorder-free system with sub-linear evolution-operator OSEE.
\S\ref{sec:mechanism} then derives the mechanism. The key quantity is not
integrability but the \textbf{number of defects that can cross the cut}: if the
sector containing the cut is spanned by configurations of $p$ domain walls with
$p$ bounded, then $\chi(t) = O(t^{\,p})$, and if $p$ is extensive $\chi$ grows
exponentially. The kink module of $W_{156}/W_{198}$ has $p=1$, giving
$\chi(t) = O(t)$; the hard-core module of $W_{201}/W_{108}$ has extensive $p$,
giving exponential growth. Fragmentation into modules of polynomial dimension is
therefore a disorder-free mechanism for the localization signature.

\textbf{New in rev.\ 2.} \S\ref{sec:tight} closes the gap left open in rev.\ 1.
The upper bound is sharpened to an exact combinatorial statement --- a
\emph{last-passage identity} valid for any operator and any splitting, giving
$\operatorname{rank}(P_RA^tP_L)\le t\cdot\operatorname{rank}(P_RAP_L)$ --- and
the cross transfer of the kink walk is shown to be \emph{rank one}, in a cell
basis in which the walk is a two-band chiral quantum walk. A matching lower
bound follows from an explicit rank-$t$ factorisation whose two vector families
are independent by causality. The result is the two-sided bound
$2t \le \chi_{\rm module}(t)\le 2t+2$, so $\chi=\Theta(t)$, with the measured
value exactly $2t+1$ for both rules, every chain length and every cut tested.
For the \emph{full} chain the exact law is $\chi(t)=6t-7$ for
$4\le t\le N/2$, verified at $N=8,10,12,14,16$; the slope $6$ itself is recorded
as measured, together with three ruled-out derivations of it.
\end{abstract}

\tableofcontents

\section{Protocol: the literature, sorted by relevance}
\label{sec:protocol}

Two objects must be kept apart throughout, because the literature treats them
very differently and conflating them is the trap this project already fell into
once:
\begin{itemize}
\item \textbf{(L)} the operator entanglement of a \emph{local} operator $O(t) =
U^\dagger O U$ in the Heisenberg picture;
\item \textbf{(U)} the operator entanglement --- equivalently the Schmidt rank
of the Choi vector, equivalently the MPO bond dimension --- of the
\emph{evolution operator} $U(t)$ itself.
\end{itemize}
Our measurement is \textbf{(U)}. Sub-linear results for (L) do not bear on it.

\subsection{Tier 1: directly on the claim}

\begin{longtable}{p{0.28\textwidth}p{0.10\textwidth}p{0.54\textwidth}}
\toprule
reference & object & finding that bears on our claim \\
\midrule
\endhead
Dubail, J.\ Phys.\ A \textbf{50}, 234001 (2017) & \textbf{(U)} &
  \textbf{The decisive baseline.} ``The OSEE of the evolution operator
  $U(t)=e^{-iHt}$ increases linearly with $t$, unless the Hamiltonian is in a
  localized phase.'' Explicitly contrasts this with local operators, whose OSEE
  ``grows sublinearly (perhaps logarithmically)''. Linear OSEE $\Rightarrow$
  $\chi$ exponential. \emph{Localization is named as the only exception.} \\
\addlinespace
Zhou \& Luitz, PRB \textbf{95}, 094206 (2017) & \textbf{(U)} &
  Classifies the growth of operator EE \emph{of the time evolution operator}
  across chaotic, MBL and Floquet systems: linear, power-law and logarithmic
  growth are all observed, before extensive saturation. Crucially the two
  sub-linear classes are the \emph{random-field} Heisenberg model --- power law
  at weak disorder, logarithmic at strong disorder. The most chaotic Floquet
  model saturates at the Page value. \\
\addlinespace
Thakur, Pain \& Roy, PRB \textbf{111}, 174206 (2025) & \textbf{(U)} &
  Operator entanglement of the time-evolution operator in the MBL phase, showing
  a \emph{logarithmic light cone}, derived from eigenstate correlations without
  invoking $\ell$-bits. This is the closest published phenomenology to ours ---
  same object, same functional form --- but obtained \emph{from disorder}. \\
\addlinespace
Prosen \& \v{Z}nidari\v{c}, PRE \textbf{75}, 015202 (2007) & (L)/state &
  The origin of the exponential-versus-polynomial dichotomy for MPO bond
  dimension: $D_\epsilon(t)\propto e^{ct}$ in quantum chaotic systems,
  polynomial when integrable. \emph{But the evolved objects are observables and
  thermal states, not the propagator}, so this does not license
  ``integrable $\Rightarrow$ polynomial $\chi$'' for $U(t)$. \\
\bottomrule
\end{longtable}

\subsection{Tier 2: the disorder-free sub-linear results, all for local operators}

\begin{longtable}{p{0.28\textwidth}p{0.10\textwidth}p{0.54\textwidth}}
\toprule
reference & object & finding \\
\midrule
\endhead
Alba, Dubail \& Medenjak, PRL \textbf{122}, 250603 (2019) & (L) &
  Logarithmic upper bound on operator entanglement for \emph{all local
  operators} in the Rule~54 chain, and evidence that log growth is generic for
  interacting integrable systems. A cellular automaton, so tempting --- but the
  wrong object. \\
\addlinespace
Tan \& Prosen, arXiv:2603.19363 (2026) & (L) &
  ``Logarithmic growth of operator entanglement in a \emph{clean non-integrable}
  circuit''. Semi-ergodic dual-unitary brickwork; mechanism is that the
  Heisenberg evolution of a single traceless single-site operator lies in a
  restricted subspace. \textbf{The nearest disorder-free precedent in spirit}
  --- log growth from a restricted subspace rather than from integrability or
  disorder --- but again for a local operator, not $U(t)$. \\
\addlinespace
Prosen \& Pi\v{z}orn, PRA \textbf{76}, 032316 (2007) & (L) &
  Transverse Ising chain: OSEE of local operators \emph{saturates} for finite
  Majorana index and grows \emph{logarithmically} for infinite index. Free
  fermions; local operators. \\
\addlinespace
Pal \& Lakshminarayan, PRB \textbf{98}, 174304 (2018) & \textbf{(U)} &
  Entangling power of time-evolution operators, integrable versus
  non-integrable. The integrable kicked transverse-field Ising chain still shows
  \emph{ballistic} growth of operator von Neumann entropy; its entangling power
  saturates lower than the non-integrable case but the growth is not sub-linear.
  \emph{Integrability alone does not buy sub-linear $\chi(t)$.} \\
\bottomrule
\end{longtable}

\subsection{Tier 3: framework, and the closest methodological neighbour}

\begin{longtable}{p{0.28\textwidth}p{0.10\textwidth}p{0.54\textwidth}}
\toprule
reference & object & finding \\
\midrule
\endhead
Merbis \& Bakker, ICCS 2024, arXiv:2406.04895 & \textbf{(U)}, classical &
  \textbf{The closest method}: builds the MPO of a cellular automaton's
  transition matrix and classifies elementary CA by operator-entanglement
  growth. Types I (converges to a constant, e.g.\ their rules 136, 108), II
  (rises then falls, peak at $t\sim L$), III (linear growth to an $L$-dependent
  plateau, chaotic-like). Generic bond dimension quoted as $4^t$, i.e.\
  $S(t)=2t$. \emph{No logarithmic class appears.} Their object is the classical
  transition matrix of a deterministic CA, not a quantum evolution operator with
  Hadamards, so the rule numbers are not comparable to ours despite the
  coincidence. \\
\addlinespace
Zanardi, Zalka \& Faoro, PRA \textbf{62}, 030301 (2000); Zanardi, PRA
  \textbf{63}, 040304 (2001) & \textbf{(U)} &
  Where the object comes from: entanglement of a unitary regarded as a vector in
  Hilbert--Schmidt space, and the entangling power. Definitional; sets the
  measure we and everyone else use. \\
\addlinespace
Jonay, Huse \& Nahum, arXiv:1803.00089 (2018) & (L)/(U) &
  The membrane / line-tension coarse-graining. Predicts \emph{linear} growth for
  generic chaotic dynamics with rate set by the surface tension
  $\mathcal E(v)$. Our result sits outside this picture, which is the point. \\
\addlinespace
Hahn, Luitz \& Chalker, PRX \textbf{14}, 031029 (2024) & \textbf{(U)} &
  Eigenstate correlations of the time-evolution operator beyond ETH; uses
  Floquet models without conserved densities to isolate scrambling. Provides the
  eigenstate-correlation language in which Thakur--Pain--Roy then derive the MBL
  log light cone. \\
\addlinespace
Bandyopadhyay \& Lakshminarayan, arXiv:quant-ph/0504052 (2005) & \textbf{(U)} &
  Entangling power of quantum chaotic evolutions; early evidence that chaotic
  evolutions approach the random-matrix value. \\
\addlinespace
Musz, Ku\'s \& \.Zyczkowski, PRA \textbf{87}, 022111 (2013) & \textbf{(U)} &
  Geometry of unitary gates, perfect entanglers, unistochastic maps. Fixes what
  ``maximal'' means for the measure. \\
\addlinespace
arXiv:2603.05656 (2026), \emph{Classical simulability from operator entanglement
scaling} & (L) &
  Simulability classification by the scaling exponent $\alpha$: $\alpha\ge1$
  blocks efficient MPO approximation, $\alpha<1$ permits it over restricted
  state ensembles. Local-operator entanglement. \\
\bottomrule
\end{longtable}

\subsection{Verdict of the pass}

\begin{enumerate}
\item For the \emph{evolution operator}, the established law is linear OSEE,
i.e.\ \textbf{exponential} $\chi(t)$, with exactly one documented exception:
\textbf{localization}, which requires disorder
(Dubail; Zhou--Luitz; Thakur--Pain--Roy).
\item Integrability is \emph{not} an exception for this object. Prosen--%
\v{Z}nidari\v{c}'s polynomial result is for observables and states; Pal--%
Lakshminarayan find ballistic operator-entropy growth even for an integrable
kicked Ising chain.
\item Every disorder-free sub-linear result we found --- Alba--Dubail--Medenjak,
Prosen--Pi\v{z}orn, Tan--Prosen --- is for \emph{local operators}, which Dubail
explicitly separates from the propagator.
\item We found \textbf{no} report of sub-linear evolution-operator Schmidt-rank
growth in a clean, disorder-free system. A linear $\chi(t)$ for $U(t)$ in a
translation-invariant Floquet circuit is, on this evidence, unoccupied ground.
\end{enumerate}

\begin{remark}
The nearest conceptual precedent is Tan--Prosen: log growth in a clean
non-integrable circuit because the evolved operator is confined to a
\emph{restricted subspace}. That is the same shape of mechanism as ours ---
confinement to a small module rather than integrability or disorder --- applied
to the other object. It should be cited as the closest relative and distinguished
on the object.
\end{remark}

\section{The mechanism: defect counting bounds the Schmidt rank}
\label{sec:mechanism}

We now give the argument that a module of linear dimension forces linear
$\chi(t)$. The statement turns out not to be about dimension as such, but about
\emph{how many defects can cross the cut}.

\subsection{Setup}

Let $U$ be one period of a brick-wall circuit with strict light-cone velocity
$v$ (sites per period). Fix a cut at bond $x$ and let $\chi_x(t)$ be the
operator Schmidt rank of $U(t) = U^t$ across it --- equivalently the rank of the
Choi vector $|U(t)\rangle\rangle$ under the bipartition, equivalently the MPO
bond dimension at that bond.

The generic bound is the light cone: $U(t)$ acts as the identity outside a
causal diamond of width $2vt$ about the cut, so
$\chi_x(t) \le 4^{\,2vt}$. Everything in \S\ref{sec:protocol} that grows
exponentially is saturating a bound of this form. To do better one needs the
dynamics to explore only a small part of the operator space on the diamond.

\subsection{The kink chain}

Take $W_{156}$. Its sectors are wall-free gaps, and inside a gap the state space
is spanned by the \emph{kink states}
\[
  |k\rangle \;=\; \underbrace{1\cdots1}_{k}\,\underbrace{0\cdots0}_{l-k},
  \qquad k = 0,\dots,l,
\]
of dimension $l+1$ --- one degree of freedom, the position of a single domain
wall. Restricted to the gap,
\[
  U(t)\big|_{\rm gap} \;=\; \sum_{k,k'} A_{kk'}(t)\, |k\rangle\langle k'| ,
\]
a \emph{single-particle propagator}. Two facts make the Schmidt rank collapse.

\paragraph{Every kink state is a product across the cut.} Cutting at $x$,
\[
  |k\rangle = |L_k\rangle \otimes |R_k\rangle,
  \qquad
  \begin{cases}
    L_k = 1^k0^{\,x-k}, \; R_k = 0^{\,l-x} & k \le x,\\
    L_k = 1^{x}, \; R_k = 1^{\,k-x}0^{\,l-k} & k > x .
  \end{cases}
\]
So each term $|k\rangle\langle k'|$ is already a product operator
$\ell_{kk'}\otimes r_{kk'}$ with
$\ell_{kk'} = |L_k\rangle\langle L_{k'}|$ and
$r_{kk'} = |R_k\rangle\langle R_{k'}|$. No entanglement is generated by the
basis itself.

\paragraph{Three of the four index sectors are rank one.} Split the sum by the
side of the cut:
\begin{itemize}
\item $k,k' \le x$: then $R_k = R_{k'} = 0^{\,l-x}$, so \emph{every} term shares
the same right factor. All of these together contribute Schmidt rank $\le 1$.
\item $k,k' > x$: then $L_k = L_{k'} = 1^{x}$, the same left factor. Rank
$\le 1$.
\item $k \le x < k'$ and $k' \le x < k$: the crossing sectors. Here the left
factor is $|L_k\rangle\langle 1^x|$, labelled by $k$ alone, and the right factor
is $|0^{\,l-x}\rangle\langle R_{k'}|$, labelled by $k'$ alone. The contribution
is the rank of the coefficient block $A_{kk'}$ restricted to that range.
\end{itemize}
Causality restricts the crossing blocks: $A_{kk'}(t) = 0$ unless
$|k-k'| \le 2vt$, and a term contributes across the cut only if the kink is
within the light cone of $x$, i.e.\ $|k-x|, |k'-x| \lesssim vt$. Each crossing
block is therefore at most $(vt)\times(vt)$, of rank at most $vt$. Hence
\begin{equation}
  \boxed{\;\chi_x(t) \;\le\; 2 + 2\,vt \;=\; O(t).\;}
  \label{eq:linear}
\end{equation}

\subsection{The general statement}

Nothing above used $\delta$, the Hadamard, or integrability. What was used is
that the sector is labelled by the position of \emph{one} defect.

\begin{proposition}[Defect counting]
\label{prop:defect}
Suppose $U$ preserves a sector whose basis states are labelled by the positions
of $p$ defects on the chain, each defect state being a product state across any
cut, and suppose $U$ has strict light-cone velocity $v$. Then the operator
Schmidt rank of $U(t)$ across a cut obeys
\[
  \chi_x(t) \;=\; O\big((vt)^{\,p}\big).
\]
If instead the number of defects is extensive in the region size, the bound
degrades to the light-cone bound and $\chi_x(t)$ may grow exponentially.
\end{proposition}

\begin{proof}[Sketch]
As above: a configuration is fixed by the $p$ defect positions, so a basis
operator $|c\rangle\langle c'|$ factorises across the cut with left factor
determined by the defect positions to the left and right factor by those to the
right. Terms in which no defect lies within the light cone of the cut share a
common factor on one side and contribute rank $1$. Terms with $j \le p$ defects
inside the light cone are labelled by at most $(2vt)^{j}$ positions on each
side, so the corresponding coefficient block has rank at most $(2vt)^{j}$.
Summing over $j \le p$ gives $O((vt)^p)$.
\end{proof}

\subsection{The two families}

\begin{center}
\begin{tabular}{lllll}
\toprule
rules & gap module & dimension & defects $p$ & predicted $\chi(t)$ \\
\midrule
$W_{156}, W_{198}$ & kink chain $1^a0^b$ & $l+1$, linear & $1$ &
  $O(t)$, \textbf{linear} \\
$W_{201}, W_{108}$ & hard-core chain & $\sim\varphi^{\,l}$, exponential &
  extensive & exponential \\
\bottomrule
\end{tabular}
\end{center}

\noindent This reproduces both measured behaviours: linear $\chi(t)$ for rule
$156$, exponential for $201/108$. It also explains why the split is so sharp
despite the two families differing by a single letter of the gate word: the
$V$-context decides the wall word, the wall word decides whether a gap admits an
ascent, and that decides whether the gap carries one defect or extensively many.

\begin{remark}[What the argument does and does not give]
It gives linearity, and the identification of the exponent with the defect
number. It does \emph{not} give the measured slope $6$. It is an upper bound, so
it does not by itself prove the growth is not slower --- \S\ref{sec:tight}
supplies the matching lower bound. And it assumes the sector basis is product
across the cut, which holds for kink and hard-core configurations but must be
checked for any other module.
\end{remark}

\begin{remark}[Why this is the localization signature without localization]
In an MBL system the log light cone comes from an emergent set of quasi-local
integrals of motion which confine information transport. Here the confinement is
kinematic: the sector simply has too few configurations within the light cone
for the Schmidt rank to grow. Fragmentation supplies, without disorder, the
thing localization supplies with it --- a bound on how many distinguishable
histories can cross a cut in time $t$. That is the sense in which
Proposition~\ref{prop:defect} is the mechanism the literature pass says is
missing.
\end{remark}

\section{The tight bound}
\label{sec:tight}

Proposition~\ref{prop:defect} gives $\chi=O(t)$. This section replaces it, for
the kink module, by a two-sided bound with an exactly matching measurement. All
numbers below come from \fn{analytics/kink\_schmidt.py}, which rebuilds the
propagator from the project's exact $\mathbb Z[1/\sqrt2]$ branch machinery, not
from any earlier tabulation. One bookkeeping point: at obc0 the maximal one-wall
Krylov sector of an $N$-site chain is $\{|k\rangle = 1^k0^{N-k}: k=1,\dots,N\}$,
of dimension $l = N$ --- the all-zero configuration is frozen and sits outside
it --- so $l$ and $N$ coincide throughout this section. Closure is exact: the
amplitude leaking out of the module is $0.0$, not merely small.

\subsection{Step 1: the reshuffle collapses onto the coefficient matrix}

Fix the cut at bond $x$ and recall from \S\ref{sec:mechanism} that
$|k\rangle = |L_{a}\rangle\otimes|R_{b}\rangle$ with $a = \min(k,x)$,
$b = \max(k,x)$. The pair $(a,b)$ determines $k$, so the operator Schmidt
decomposition of $U(t)|_{\rm module}$ is nothing but a re-indexing of the
coefficient matrix $A^t$: rows $(a,a')$, columns $(b,b')$. The four index
sectors land as follows.

\begin{center}
\begin{tabular}{llll}
\toprule
index sector & row & column & content \\
\midrule
$k,k'\le x$ & $(k,k')$ & $(x,x)$ & the whole $LL$ block, in one column \\
$k,k'> x$ & $(x,x)$ & $(k,k')$ & the whole $RR$ block, in one row \\
$k\le x<k'$ & $(k,x)$ & $(x,k')$ & the crossing block $C_{LR}$ \\
$k'\le x<k$ & $(x,k')$ & $(k,x)$ & the crossing block $C_{RL}$ \\
\bottomrule
\end{tabular}
\end{center}

\noindent The two crossing blocks occupy disjoint rows and disjoint columns, and
the remaining content is confined to the single row $(x,x)$ and the single
column $(x,x)$. Hence, with no approximation,
\begin{equation}
  \operatorname{rank}C_{LR} + \operatorname{rank}C_{RL}
  \;\le\; \chi_x(t) \;\le\;
  2 + \operatorname{rank}C_{LR} + \operatorname{rank}C_{RL}.
  \label{eq:sandwich}
\end{equation}
The whole problem is now the rank of one crossing block.

\subsection{Step 2: a last-passage identity, and the transfer rank}

\begin{proposition}[Transfer rank]
\label{prop:transfer}
Let $A$ act on a finite-dimensional space and let $P_L + P_R = \mathbb 1$ be
complementary projectors. Then for every $t\ge1$
\begin{equation}
  P_R A^t P_L \;=\; \sum_{s=0}^{t-1}
    \big(P_R A P_R\big)^{t-1-s}\,\big(P_R A P_L\big)\,\big(P_L A^{s} P_L\big),
  \label{eq:lastpass}
\end{equation}
and consequently
\[
  \operatorname{rank}\big(P_R A^t P_L\big)
  \;\le\; t\cdot\operatorname{rank}\big(P_R A P_L\big).
\]
\end{proposition}

\begin{proof}
Expand $A^t$ as a sum over sequences of $L/R$ labels, one per time step. Every
sequence contributing to $P_RA^tP_L$ starts in $L$ and ends in $R$, so it has a
unique \emph{last} time $s$ at which the label is $L$, with $0\le s\le t-1$.
Grouping the sequences by that $s$ gives exactly \eqref{eq:lastpass}: the factor
$P_LA^sP_L$ collects all histories up to time $s$ (unrestricted in between), the
factor $P_RAP_L$ is the step out of $L$ at time $s+1$, and
$(P_RAP_R)^{t-1-s}$ the remaining steps, which stay in $R$ by construction. Each
summand contains $P_RAP_L$ as a factor, so each has rank at most
$\operatorname{rank}(P_RAP_L)$, and rank is subadditive.
\end{proof}

The identity is checked numerically to $\le 2.2\times10^{-16}$ for the kink
propagator at $N=24,30$ and $t\le10$, and for a random $10\times10$ unitary; it
is elementary, but it is the step that converts a \emph{static}, one-period
quantity into a linear-in-$t$ bound.

It remains to compute $\operatorname{rank}(P_RAP_L)$ for the kink walk. Group
the kink positions into cells $\{2j-1,2j\}$ and pass to
\[
  |p_j\rangle = \tfrac1{\sqrt2}\big(|2j-1\rangle + |2j\rangle\big),
  \qquad
  |m_j\rangle = \tfrac1{\sqrt2}\big(|2j-1\rangle - |2j\rangle\big).
\]
In this basis the exact one-period propagator of $W_{156}$ reads, in the bulk,
\begin{align}
  A\,|p_j\rangle &= \tfrac12\big(|m_{j-1}\rangle + |m_j\rangle
    - |p_j\rangle + |p_{j+1}\rangle\big), \nonumber\\
  A\,|m_j\rangle &= \tfrac12\big(|m_{j-1}\rangle - |m_j\rangle
    - |p_j\rangle - |p_{j+1}\rangle\big),
  \label{eq:cellwalk}
\end{align}
i.e.\ a two-band quantum walk hopping by at most one cell per period. Its
inter-cell transfer operators are
\begin{equation}
  T_+ \;=\; \tfrac12\,|p\rangle\big(\langle p| - \langle m|\big),
  \qquad
  T_- \;=\; \tfrac12\,|m\rangle\big(\langle p| + \langle m|\big),
  \label{eq:transfer}
\end{equation}
\emph{both of rank one} ($W_{198}$ is the mirror of $W_{156}$ and reflection
sends $k\mapsto l+1-k$, so its cell tiling is the shifted pairing
$\{2j,2j+1\}$ and the roles of $p$ and $m$ are exchanged).
The coin state $p$ is the sole right-moving channel and $m$ the sole left-moving
one: the walk is chiral, with one channel per chirality per cell. Measured
directly, $\operatorname{rank}(P_RAP_L) = 1$ at \emph{every} cut
$x=1,\dots,l-1$, for both rules at $N=12,20,24$, so
Proposition~\ref{prop:transfer} gives
$\operatorname{rank}C_{LR}\le t$, likewise for $C_{RL}$.

\subsection{Step 3: the matching lower bound}

Because $P_RAP_L = ab^{\mathsf T}$ is rank one, \eqref{eq:lastpass} is an
explicit rank-$t$ factorisation:
\[
  P_R A^t P_L \;=\; \sum_{s=1}^{t} u_s v_s^{\mathsf T},
  \qquad
  u_s = (P_RAP_R)^{t-s} a, \qquad
  v_s = \big(P_L A^{s-1} P_L\big)^{\mathsf T} b .
\]
If $\{u_s\}_{s\le t}$ and $\{v_s\}_{s\le t}$ are each linearly independent, the
factorisation is $UV^{\mathsf T}$ with both factors of full column rank $t$, so
the rank is exactly $t$. Independence follows from the light cone: by
\eqref{eq:transfer} the rightward amplitude per period is exactly $\tfrac12$ and
never vanishes, so $u_s$ reaches exactly $t-s$ cells to the right of the cut
with leading amplitude $2^{-(t-s)}$, and $v_s$ exactly $s-1$ cells to the left.
The reaches are pairwise distinct, hence each family is independent. Measured:
$\dim\operatorname{span}\{u_s\} = \dim\operatorname{span}\{v_s\}
= \operatorname{rank} = t$ for all $t\le 9$ at $N=24,30$, both rules, with the
factorisation reproducing $P_RA^tP_L$ to $\le2.2\times10^{-16}$.

\subsection{The law}

Combining \eqref{eq:sandwich} with $\operatorname{rank}C_{LR}
= \operatorname{rank}C_{RL} = t$:

\begin{theorem}[Kink module]
\label{thm:tight}
For $W_{156}$ and $W_{198}$ on their maximal one-wall module of dimension $l$,
and any cut $x$ with $1\le t<\min(x,\,l-x)$ --- i.e.\ as long as the light cone
has not reached either end of the chain ---
\[
  2t \;\le\; \chi_x(t) \;\le\; 2t+2 ,
  \qquad\text{so}\qquad \chi_x(t) = \Theta(t).
\]
\end{theorem}

\noindent The measured value is the middle one, and the saturation is a clean
truncation:
\begin{equation}
  \chi_x(t) \;=\; \min\big(2t+1,\; \chi_x^{\rm sat}\big),
  \qquad t\ge1,
  \label{eq:exactmodule}
\end{equation}
with $\chi_x^{\rm sat}$ a $t$-independent, cut-dependent ceiling. Verified
exhaustively for both rules at $N=12,16,20$, \emph{every} cut $x=1,\dots,N-1$
and $t$ up to $3N$, and at $N=24,30$ for $x=3,\,N/2,\,N/2+1$. At $t=0$ the value
is $2$: the identity on the module is not a product operator across the cut,
being $\mathbb 1_{\rm left}\otimes P_R + P_L\otimes\mathbb 1_{\rm right}$. This
is the statement rev.\ 1 left open. The ceiling $\chi_x^{\rm sat}$ is a
finite-chain edge effect, small when the cut is near an end; no closed form for
it is claimed here.

\subsection{The full chain}

On the full $2^N$ space the exact operator Schmidt rank of $U^t$ across the
middle bond is
\begin{equation}
  \boxed{\;\chi(t) \;=\; 6t-7,\qquad 4\le t\le N/2.\;}
  \label{eq:full}
\end{equation}
Measured series (identical for $W_{156}$ and $W_{198}$, obc0, $x=N/2$):

\begin{center}
\begin{tabular}{llll}
\toprule
$N$ & $\chi(t)$, $t=0,1,2,\dots$ & law holds for & first deviation \\
\midrule
$8$  & $1,4,8,12,17,20,21,\dots$ & $t\le4$ & $t=5$ \\
$10$ & $1,4,8,12,17,23,27,30,31,\dots$ & $t\le5$ & $t=6$ \\
$12$ & $1,4,8,12,17,23,29,33,37,41,\dots$ & $t\le6$ & $t=7$ \\
$14$ & $1,4,8,12,17,23,29,35,39,43$ & $t\le7$ & $t=8$ \\
$16$ & $1,4,8,12,17,23,29,35,41,45,49$ & $t\le8$ & $t=9$ \\
\bottomrule
\end{tabular}
\end{center}

\noindent The increments are $3,4,4,5$ and then a constant $6$; the deviations
past $t=N/2$ are the finite-chain saturation, and the window grows by one period
per two sites, exactly as a light cone should. The $N=14$ and $N=16$ rows were
computed with a sparse propagator and a randomised rank estimator (oversampled
to $160$ probe vectors, so the reported rank is exact whenever it is below the
cap, which it always is).

\begin{remark}[The slope $6$ is not yet derived --- and three routes are closed]
Theorem~\ref{thm:tight} explains the linearity and gives slope $2$ per
straddling gap. The factor $3$ relating that to \eqref{eq:full} is not derived
here, and the obvious derivations do not work:
\begin{enumerate}
\item \emph{Not a sum over Krylov sectors.} The reshuffle is supported on entry
$(y,y')$ only when $y,y'$ share a sector, so it splits into blocks of disjoint
support; but the left Schmidt vectors are shared \emph{between} sectors, and the
sum is wildly larger than the total. At $N=12$ (377 sectors)
$\sum_S\chi_S(t) = 473,569,700,747,763,749,731,\dots$ against
$\chi(t)=1,4,8,12,17,23,29,\dots$ --- not even monotone. Fragmentation alone does
not make the MPO small; the sharing does.
\item \emph{Not a sum over growing sectors either.} Only $12$ of the $377$
sectors have a $\chi_S$ still growing at $t=6$, each with slope $2$ while
unsaturated, for a total of $24$ where the full chain increments by $6$.
\item \emph{Not an accumulating channel basis.} The left Schmidt subspace does
not grow as a filtration: projecting the $t{+}1$ subspace against the $t$ one
leaves $20$--$21$ new directions out of $23$--$29$, so the subspace rotates
almost completely each period and there is no fixed set of ``six new channels''
to count. Every Schmidt vector also acts non-trivially out to the far end of the
half-chain, so the new directions cannot be localised at the light-cone edge.
\end{enumerate}
The slope is therefore a genuinely global count over the wall grammar, and
deriving it needs the MPO of $U^t$ itself rather than any sector-wise argument.
\end{remark}

\section{What would settle it}

\begin{enumerate}
\item \textbf{The slope $6$.} The one quantitative gap left. \eqref{eq:full} is
exact over the whole accessible window and the three cheap explanations above
are excluded, so this needs the transfer-matrix structure of the MPO of $U^t$.
\item \textbf{$\chi(t)$ for $W_{201}, W_{108}$.} Held by the experimental side
of the project; the linear-versus-exponential contrast is the whole claim and
Proposition~\ref{prop:defect} predicts exponential growth there because the
hard-core module has extensive defect number.
\item \textbf{Beyond one straddling gap.} Theorem~\ref{thm:tight} is the $p=1$
case. Confirming the $O(t^{\,p})$ prediction of
Proposition~\ref{prop:defect} on a module with $p=2$ would turn defect counting
from an explanation of two data points into a law.
\end{enumerate}

\begin{thebibliography}{99}
\bibitem{Dubail2017} J.~Dubail,
  \emph{Entanglement scaling of operators: a conformal field theory approach,
  with a glimpse of simulability of long-time dynamics in $1+1$d},
  J.\ Phys.\ A \textbf{50}, 234001 (2017); arXiv:1612.08630.
\bibitem{ZhouLuitz2017} T.~Zhou and D.~J.~Luitz,
  \emph{Operator entanglement entropy of the time evolution operator in chaotic
  systems}, Phys.\ Rev.\ B \textbf{95}, 094206 (2017); arXiv:1612.07327.
\bibitem{Thakur2025} R.~Thakur, B.~Pain and S.~Roy,
  \emph{Logarithmic entanglement light cone from eigenstate correlations in the
  many-body localized phase}, Phys.\ Rev.\ B \textbf{111}, 174206 (2025);
  arXiv:2504.02815.
\bibitem{PZ2007} T.~Prosen and M.~\v{Z}nidari\v{c},
  \emph{Is the efficiency of classical simulations of quantum dynamics related
  to integrability?}, Phys.\ Rev.\ E \textbf{75}, 015202 (2007).
\bibitem{PP2007} T.~Prosen and I.~Pi\v{z}orn,
  \emph{Operator space entanglement entropy in a transverse Ising chain},
  Phys.\ Rev.\ A \textbf{76}, 032316 (2007); arXiv:0706.2480.
\bibitem{ADM2019} V.~Alba, J.~Dubail and M.~Medenjak,
  \emph{Operator entanglement in interacting integrable quantum systems: the
  case of the Rule 54 chain},
  Phys.\ Rev.\ Lett.\ \textbf{122}, 250603 (2019); arXiv:1901.04521.
\bibitem{TanProsen2026} M.~T.~Tan and T.~Prosen,
  \emph{Logarithmic growth of operator entanglement in a clean non-integrable
  circuit}, arXiv:2603.19363 (2026).
\bibitem{PalLak2018} R.~Pal and A.~Lakshminarayan,
  \emph{Entangling power of time-evolution operators in integrable and
  nonintegrable many-body systems},
  Phys.\ Rev.\ B \textbf{98}, 174304 (2018); arXiv:1805.11632.
\bibitem{Merbis2024} W.~Merbis and C.~C.~Bakker,
  \emph{Operator entanglement growth quantifies complexity of cellular
  automata}, ICCS 2024; arXiv:2406.04895.
\bibitem{ZZF2000} P.~Zanardi, C.~Zalka and L.~Faoro,
  \emph{Entangling power of quantum evolutions},
  Phys.\ Rev.\ A \textbf{62}, 030301 (2000).
\bibitem{Zanardi2001} P.~Zanardi,
  \emph{Entanglement of quantum evolutions},
  Phys.\ Rev.\ A \textbf{63}, 040304 (2001); arXiv:quant-ph/0010074.
\bibitem{JHN2018} C.~Jonay, D.~A.~Huse and A.~Nahum,
  \emph{Coarse-grained dynamics of operator and state entanglement},
  arXiv:1803.00089 (2018).
\bibitem{HLC2024} D.~Hahn, D.~J.~Luitz and J.~T.~Chalker,
  \emph{Eigenstate correlations, the eigenstate thermalization hypothesis, and
  quantum information dynamics in chaotic many-body quantum systems},
  Phys.\ Rev.\ X \textbf{14}, 031029 (2024); arXiv:2309.12982.
\bibitem{BL2005} J.~N.~Bandyopadhyay and A.~Lakshminarayan,
  \emph{Entangling power of quantum chaotic evolutions via operator
  entanglement}, arXiv:quant-ph/0504052 (2005).
\bibitem{MKZ2013} M.~Musz, M.~Ku\'s and K.~\.Zyczkowski,
  \emph{Unitary quantum gates, perfect entanglers, and unistochastic maps},
  Phys.\ Rev.\ A \textbf{87}, 022111 (2013).
\bibitem{Sim2026} \emph{Classical simulability from operator entanglement
  scaling}, arXiv:2603.05656 (2026).
\end{thebibliography}

\end{document}
